% ---------------------------------------------------------------------
\section{Related Work}\label{sec:related}
% ---------------------------------------------------------------------

\subsection{Reinforcement learning for EV charging and V2G}\label{sec:rw-rl}

Deep RL for EV charging has moved from single-vehicle agents \cite{li2020constrained} to stations with photovoltaics, storage and V2G \cite{korkas2024distributed,zhang2025dumes}, to coordinated stations \cite{song2025longterm} and to fleets under distribution-grid limits \cite{zhang2024hamlet,caballe2026largescale}.
Two reviews list safety, dependence on simulation and the lack of a standard benchmark among the open problems \cite{qiu2023reinforcement,michailidis2025reinforcement}.
EV2Gym models a CPO whose chargers sit under transformers with inflexible load and solar generation \cite{orfanoudakis2025ev2gym}.
Its V2G reward penalises transformer overload and low user satisfaction, and in its V2G benchmark penalty-based agents earned profit but left many users short of their requested energy.
The Adaptive Charging Network (ACN) solves a convex program in receding horizon with hard infrastructure limits and best-effort energy delivery \cite{lee2021adaptive}.
An economic model predictive control (MPC) inside EV2Gym enforces hard departure and transformer constraints \cite{diazlondono2026open}.

\subsection{Constraint handling and safe reinforcement learning}\label{sec:rw-safe}

Weighted violation terms in the reward do not prevent violations at test time \cite{tabas2022computationally,orfanoudakis2025ev2gym}.
The constrained Markov decision process (CMDP) bounds an expected cumulative cost instead \cite{altman2021constrained}, solved by Lagrangian \cite{tessler2018reward} or trust-region \cite{achiam2017constrained} updates.
Reviews note that such bounds hold for an expected cumulative cost rather than at each step \cite{gu2024review} and classify Lagrangian methods as soft guarantees \cite{su2025review}.
CMDP formulations have been applied to EV charging with voltage limits \cite{fan2024safety}, transformer thermal stress \cite{zhang2025modelfree,zhang2026safe} and V2G aggregators \cite{eland2026safe}.
A Lagrangian multi-agent method also encodes the recent history of each station with a transformer \cite{huang2026safe}.

A second family corrects each action before execution, through shields \cite{alshiekh2018safe}, learned linear safety layers \cite{dalal2018safe} or quadratic programs with a fallback controller \cite{ceusters2023adaptive}.
EV examples include rule-based shields \cite{guan2024rule}, action-replacement layers \cite{zhang2025dumes}, grid-code shielding \cite{demirel2025multiobjective} and power-flow shields \cite{zaman2026graph}.
Such a filter is only as safe as its internal model, and the authors of the HetGPS preprint make no pointwise feasibility claim for its learned risk model with a linearised power-flow projection \cite{wang2026hetgps}.
Members of the EV2Gym team have also combined a graph-enhanced twin delayed deep deterministic policy gradient (TD3) agent \cite{fujimoto2018addressing} with a safety layer \cite{orfanoudakis2026can}.
According to the chapter abstract, the layer is intended to preserve charging requirements and enforce power limits.

A third family makes the policy output feasible by construction.
Differentiable convex optimisation layers pass gradients through an exact solve \cite{amos2017optnet,agrawal2019differentiable}, and Chen et al. used such a projection inside energy-system policies \cite{chen2021enforcing}.
Gauge maps give feasible points directly \cite{li2023learning}, and Tabas and Zhang used one to prove that a learned frequency controller stays in a controlled-invariant safe set \cite{tabas2022computationally}.
For EV stations, a preprint by Wan et al. proves per-step feasibility of a quadratic-program projection that covers transformer and charger-current limits but not departure energy \cite{wan2025logicinformed}.
A preprint by Yan et al. derives per-EV bounds from remaining energy and time, and its formal result concerns first-order power-balance feasibility \cite{yan2026hierarchical}.
Projection also alters learning, because naive projection biases policy gradients \cite{gros2020safe} and an active projection has zero gradients \cite{kasaura2023benchmarking}.
For stochastic policies, a recent preprint shows that projection inside the environment is equivalent to projection as a policy layer \cite{markgraf2025safe}.

\subsection{Size-agnostic architectures}\label{sec:rw-scale}

Deep Sets \cite{zaheer2017deep}, self-attention \cite{vaswani2017attention} and the Set Transformer \cite{lee2019set} process sets of any size with shared weights, and UPDeT applied this idea to multi-agent tasks of different sizes \cite{hu2021updet}.
Size transfer is not automatic, since neural routing solvers trained on small graphs often fail on larger ones \cite{joshi2022learning}.
Early size-independent states in EV charging came from aggregation, and Sadeghianpourhamami et al. also reported the gap to a perfect-knowledge optimum under scaled arrival rates \cite{sadeghianpourhamami2020definition}.
Aggregated representations, however, hide per-vehicle detail or move it into disaggregation rules \cite{zhang2022transfer,salesortiz2024efficient}.

The preprint version of the heterogeneous graph policy (EV-GNN) coordinates up to 1000 charge points, with a separate model per scale and no constraint enforced by construction \cite{orfanoudakis2025scalable}.
In contrast, the graph-enhanced decision transformer (GNN-DT), trained on 25 chargers, was tested without retraining on 5 to 100 chargers, and its performance degraded as size grew \cite{orfanoudakis2026graph}.
A HetGPS policy trained with 1,236 vehicles transfers without retraining to feeders with 2,239 and 3,218 vehicles \cite{wang2026hetgps}.
Similarly, a graph actor for voltage control has been evaluated in simulation on a larger feeder than in training \cite{jin2026telemetry}.

In these size-agnostic controllers, constraints enter through reward terms, imitation data or heuristic filters.

\subsection{Laxity-based scheduling and departure uncertainty}\label{sec:rw-laxity}

Chen et al. proved that a schedule is feasible exactly when every laxity, the slack between the time left and the time still needed to finish charging, stays non-negative \cite{chen2022smoothed}.
Their sLLF rule clips each rate around a common level under the station power limit, and their bounds indicate that online feasibility generally requires extra power.
A preprint shows that hard deadlines and a transformer limit can become jointly infeasible under congestion \cite{reddy2026deadline}, and least-laxity-first (LLF) priorities perform well under cable and transformer limits \cite{vanhuffelen2025grid,christensen2025multiagent}.
These methods assume that users declare departure time and energy at arrival.
Plug-in durations are hard to predict \cite{dyo2022behavior}, users can gain by under-reporting deadlines \cite{ferragut2022scheduling}, and departure forecasts improve prioritisation \cite{frendo2021improving}.
A recent preprint trains a departure forecaster end-to-end with the charging policy but keeps unmet energy as a soft penalty and gives no coverage statement \cite{gabriele2026forecasting}.
Conformal prediction gives distribution-free coverage under exchangeability \cite{lei2018distributionfree,angelopoulos2023conformal} and has tightened constraints in safe planning \cite{lindemann2023safe}.

\subsection{Positioning}\label{sec:rw-position}

As seen in \Tab{tab:positioning}, each ingredient of our design already exists in some form.
Size-agnostic policies with zero-shot transfer exist for EV charging \cite{orfanoudakis2026graph,wang2026hetgps}, provable action-correction layers exist in safe RL \cite{tabas2022computationally,wan2025logicinformed}, and a closely related clip-around-a-common-level rule appears in sLLF \cite{chen2022smoothed}.
Per-EV latest-start bounds appear in the scheduler of Yan et al. \cite{yan2026hierarchical}, and train-small, test-large evaluation with a gap to an optimum dates back at least to \cite{sadeghianpourhamami2020definition}.
Three closely related recent works were available to us only as abstracts \cite{orfanoudakis2026can,song2025longterm,eland2026safe}, and we describe them only at that level.
To our knowledge, no prior work joins a size-agnostic learned policy, a per-step transformer guarantee and a multi-step departure argument in one controller.
% Author note (not printed): re-check orfanoudakis2026can, song2025longterm, eland2026safe and the published EV-GNN article in full text before submission (gap analysis R1-R4).

\begin{table}[t]
\centering
\caption{Positioning against the closest works. ``Penalty'' means a reward term; ``hard'' means enforced by construction.}
\label{tab:positioning}
\footnotesize
\setlength{\tabcolsep}{3pt}
\begin{tabular}{@{}>{\raggedright}p{2.6cm}>{\raggedright}p{2.0cm}>{\raggedright}p{2.0cm}>{\raggedright}p{2.0cm}>{\raggedright}p{2.2cm}>{\raggedright}p{1.9cm}>{\raggedright\arraybackslash}p{2.0cm}@{}}
\toprule
Work & Size-agnostic policy & Departure energy & Transformer limit & Guarantee & Zero-shot size test & Gap to optimum \\
\midrule
EV2Gym agents \cite{orfanoudakis2025ev2gym} & no (flat) & penalty & penalty & none & no & optimum as baseline \\
EV-GNN, preprint version \cite{orfanoudakis2025scalable} & yes (graph) & not by construction & not by construction & none & separate model per scale & optimisation baseline per scale \\
GNN-DT \cite{orfanoudakis2026graph} & yes (graph) & penalty & penalty & none & 25 to 5--100 chargers & near-oracle reward (journal abstract) \\
HetGPS, preprint \cite{wang2026hetgps} & yes (graph attention) & reward and heuristic & not constrained & empirical (voltage filter) & 1,236 to 3,218 vehicles & no \\
Yan et al., preprint \cite{yan2026hierarchical} & no & hard (per-EV rule) & not modelled & first-order power balance & no & gap to MPC \\
Wan et al., preprint \cite{wan2025logicinformed} & no & not modelled & hard (quadratic-program projection) & per-step feasibility & no & no \\
sLLF \cite{chen2022smoothed} & rule, any size & laxity rule & hard (common level) & feasibility theorems & not learned & no \\
Sadeghian\-pourhamami et al. \cite{sadeghianpourhamami2020definition} & yes (aggregate state) & penalty & not modelled & none & scaled arrivals & yes \\
Orfanoudakis et al. \cite{orfanoudakis2026can} (abstract) & graph policy & safety layer (stated aim) & safety layer (stated aim) & not stated in abstract & not stated in abstract & not stated in abstract \\
\midrule
This work & yes (attention and Deep Sets over charger tokens) & per-EV bound, conditional (Thm.~1b) & hard (common shift) & transformer: every step (A3, A6); departure: per EV, conditional & 10/20 to 200 chargers & soft- and hard-constrained optima \\
\bottomrule
\end{tabular}
\end{table}
